% master: main
% format: latex

The automation of tasks that are dangerous, tedious, or otherwise undesirable
for humans are one of the main benefits of intelligent and programmable robots.
Another benefit is that of automating tasks that would otherwise be unfeasible
to do with human resources. The use of robotics as a tool to enhance human
capital, as well as automate tasks that are not desirable for humans to perform,
are one of the great advantages of using said platforms.
The Aquapod is an amphibious platform design to aid in search and rescue
operations as well as environmental monitoring. It is a relatively inexpensive
robot that can be used to carry sensors to locations on land, on the surface of
a body of water, while submerged, and any combination of the two. 
In order to carry sensors from once place to the next controllably, robustly,
and reliably, the Aquapod requires the development of an all-terrain powertrain
to match its target environment. 

%% Introduce the screwdrive? .... 

%% Existing structure?
Furthermore, the screwdrive powertrain has been designed to work with existing
Aquapod platforms without requiring changes to the existing robot in order to
function with the screwdrive. 

%% Dual modality 
The screwdrive drivetrain is a dual-modality powertrain that leverages the
existing tumbling modality of the Aquapod and Adelopod family of robots and adds
and additional mode that allows it to position itself controllably to tumble,
but, more importantly, allows the control of the robot while in a body of water
by using the screws as propellers. Furthermore, maintaining the tumbling mode,
it's possible to control the direction of the propulsive force provided by the
screws, overall enhancing the mobility of the Aquapod in a body of water and
achieve controllability over its six degrees of freedom. 
